\documentclass[10pt,a4paper]{amsart}
\usepackage{amsmath,amssymb,amsthm,graphicx}
\usepackage{verbatim}
\newenvironment{reference}{\comment}{\endcomment}
%\newenvironment{reference}{}{}
\newenvironment{slogan}{\comment}{\endcomment}
\newenvironment{history}{\comment}{\endcomment}

% For commutative diagrams we use Xy-pic
\usepackage[all]{xy}

% We use 2cell for 2-commutative diagrams.
\xyoption{2cell}
% Theorem environments.
%
\theoremstyle{plain}
\newtheorem{theorem}[subsection]{Theorem}
\newtheorem{proposition}[subsection]{Proposition}
\newtheorem{lemma}[subsection]{Lemma}

\theoremstyle{definition}
\newtheorem{definition}[subsection]{Definition}
\newtheorem{example}[subsection]{Example}
\newtheorem{exercise}[subsection]{Exercise}
\newtheorem{situation}[subsection]{Situation}

\theoremstyle{remark}
\newtheorem{remark}[subsection]{Remark}
\newtheorem{remarks}[subsection]{Remarks}

\numberwithin{equation}{subsection}

% Macros
%
\def\lim{\mathop{\mathrm{lim}}\nolimits}
\def\colim{\mathop{\mathrm{colim}}\nolimits}
\def\Spec{\mathop{\mathrm{Spec}}}
\def\Hom{\mathop{\mathrm{Hom}}\nolimits}
\def\Ext{\mathop{\mathrm{Ext}}\nolimits}
\def\SheafHom{\mathop{\mathcal{H}\!\mathit{om}}\nolimits}
\def\SheafExt{\mathop{\mathcal{E}\!\mathit{xt}}\nolimits}
\def\Sch{\mathit{Sch}}
\def\Mor{\mathop{\mathrm{Mor}}\nolimits}
\def\Ob{\mathop{\mathrm{Ob}}\nolimits}
\def\Sh{\mathop{\mathit{Sh}}\nolimits}
\def\NL{\mathop{N\!L}\nolimits}
\def\CH{\mathop{\mathrm{CH}}\nolimits}
\def\proetale{{pro\text{-}\acute{e}tale}}
\def\etale{{\acute{e}tale}}
\def\QCoh{\mathit{QCoh}}
\def\Ker{\mathop{\mathrm{Ker}}}
\def\Im{\mathop{\mathrm{Im}}}
\def\Coker{\mathop{\mathrm{Coker}}}
\def\Coim{\mathop{\mathrm{Coim}}}

% Boxtimes
%
\DeclareMathSymbol{\boxtimes}{\mathbin}{AMSa}{"02}

%
% Macros for moduli stacks/spaces
%
\def\QCohstack{\mathcal{QC}\!\mathit{oh}}
\def\Cohstack{\mathcal{C}\!\mathit{oh}}
\def\Spacesstack{\mathcal{S}\!\mathit{paces}}
\def\Quotfunctor{\mathrm{Quot}}
\def\Hilbfunctor{\mathrm{Hilb}}
\def\Curvesstack{\mathcal{C}\!\mathit{urves}}
\def\Polarizedstack{\mathcal{P}\!\mathit{olarized}}
\def\Complexesstack{\mathcal{C}\!\mathit{omplexes}}
% \Pic is the operator that assigns to X its picard group, usage \Pic(X)
% \Picardstack_{X/B} denotes the Picard stack of X over B
% \Picardfunctor_{X/B} denotes the Picard functor of X over B
\def\Pic{\mathop{\mathrm{Pic}}\nolimits}
\def\Picardstack{\mathcal{P}\!\mathit{ic}}
\def\Picardfunctor{\mathrm{Pic}}
\def\Deformationcategory{\mathcal{D}\!\mathit{ef}}
\usepackage{cleveref}
\crefname{originalreference}{original reference}{original references}
\Crefname{originalreference}{Original reference}{Original references}
\makeatletter
\newlabel{lemma-flat-proper-perfect-direct-image-general}{{0B91}{source}{}{sourcepointer-1}{}}
\newlabel{lemma-flat-proper-perfect-direct-image-general@cref}{{[originalreference][0][]0B91}{[1][0][]source}{}{}{}}
\makeatother
\crefname{originalreference}{original reference}{original references}
\Crefname{originalreference}{Original reference}{Original references}
\makeatletter
\newlabel{lemma-open-where-cohomology-in-degree-i-rank-r}{{0BDK}{source}{}{sourcepointer-2}{}}
\newlabel{lemma-open-where-cohomology-in-degree-i-rank-r@cref}{{[originalreference][0][]0BDK}{[1][0][]source}{}{}{}}
\makeatother
\title{Openness on the base of the locus where a base-change compatible direct image is a finite locally free module of fixed rank in a fixed degree}
\author{The Stacks Project Authors}
\date{stacks-project:spaces-perfect.tex:local-snapshot-20261012; GNU FDL 1.2 or later}
\begin{document}
\maketitle
\noindent\textit{Editorial note.} All delivered mathematical text is the author's own, copied from the pinned Stacks Project chapter file spaces-perfect.tex; no statement or proof step was rewritten, reordered or removed.\par\medskip
\noindent\textit{Editorial note.} The counted proof is the author's complete argument: Reduction-type card: the span contains the complete argument; the two cited inputs (flat proper perfect direct image with base change, and the perfect openness locus) are stated as lemmas of the same chapter.\par\medskip
\noindent\textit{Editorial note.} The label lemma-open-where-cohomology-in-degree-i-rank-r-geometric names the counted statement, and the proof environment immediately following it in the source is delivered verbatim as its proof.\par\medskip
\noindent\textit{Editorial note.} The counted statement is carried by the author's own lemma environment in the chapter source.\par\medskip
\noindent\textit{Allowed standard prerequisites (editorial).}\begin{itemize}
\item Cohomology and base change for a proper flat finitely presented morphism f : X to Y with F of finite presentation flat over Y: the object K = Rf\_* F is perfect and its formation commutes with arbitrary base change (flat proper perfect direct image lemma)
\item The perfect complex openness lemma of the same chapter: for a perfect object K of D(O\_Y) and fixed i, r there is an open V of Y such that T to Y factors through V if and only if K restricted over T is isomorphic to a finite locally free module of rank r placed in degree i (open where cohomology in degree i rank r lemma)
\end{itemize}\medskip
\ifdefined\hypertarget\hypertarget{sourcepointer-1}{}\fi
\noindent\textit{Source reference.} Original reference 0B91: sources/stacks-spaces-perfect/pinned-original-source.tex, source lines 6442--6442. Prerequisite pointer invoked at 6780; it names a standard fact of the same chapter that the counted argument uses, and every step of the counted proof built on it is retained verbatim in the delivered body.\par\medskip
\ifdefined\hypertarget\hypertarget{sourcepointer-2}{}\fi
\noindent\textit{Source reference.} Original reference 0BDK: sources/stacks-spaces-perfect/pinned-original-source.tex, source lines 6743--6743. Prerequisite pointer invoked at 6784; it names a standard fact of the same chapter that the counted argument uses, and every step of the counted proof built on it is retained verbatim in the delivered body.\par\medskip
\begin{lemma}
\label{lemma-open-where-cohomology-in-degree-i-rank-r-geometric}
Let $S$ be a scheme.
Let $f : X \to Y$ be a morphism of algebraic spaces over $S$
which is proper, flat, and of finite presentation.
Let $\mathcal{F}$ be an $\mathcal{O}_X$-module of finite presentation,
flat over $Y$. Fix $i, r \in \mathbf{Z}$.
Then there exists an open subspace
$V \subset Y$ with the following property:
A morphism $T \to Y$ factors through $V$ if and only if
$Rf_{T, *}\mathcal{F}_T$ is isomorphic to a
finite locally free module of rank $r$ placed in degree $i$.
\end{lemma}

\begin{proof}
By cohomology and base change (
Lemma \ref{lemma-flat-proper-perfect-direct-image-general})
the object $K = Rf_*\mathcal{F}$ is a perfect object of the derived
category of $Y$ whose formation commutes with arbitrary base change.
Thus this lemma follows immediately from
Lemma \ref{lemma-open-where-cohomology-in-degree-i-rank-r}.
\end{proof}

\end{document}
